% =========================================================================
% NDH-SURREALISM PREPRINT SKELETON: FROM AUTONOMIC ART TO SPECTRAL SOLAR DESIGN
% =========================================================================
\documentclass[11pt,twocolumn]{article}
\usepackage{amsmath,amssymb,amsfonts,graphicx,hyperref,microtype}

\title{\textbf{Tensor Choreography and Non-Euclidean Deposition Substrates:\\From Autonomic Art Production to Spectral Solar Engineering}}
\author{\textbf{Borealis S. Hedling} \and \textbf{The Studio Compiler Network}}
\date{November 2026}

\begin{document}
\maketitle

\begin{abstract}
This paper presents a formal framework for the deterministic translation of high-dimensional, unstable latent manifolds onto physical two-dimensional and three-dimensional substrates without geometric boundary degradation. We detail an empirical proof-of-concept pipeline (`NDH-Surrealism`) currently optimized for robotic canvas manipulation. We demonstrate that by combining low-rank spectral tensor choreography with machine-checked Coq verification enclaves, these translation invariants generalize directly to fractal architectural design and spectrally optimized solar material deposition.
\end{abstract}

% --- SECTION 1: INTRODUCTION & PROOF OF CONCEPT ---
\section{Introduction \& Empirical Baseline}
Current generative architectures operate within high-dimensional ungrounded manifolds, frequently exhibiting structural phenomena classified as hallucination geometry or metric decay. While typical machine learning pipelines employ unverified, heavy optimization routines to suppress these variations, our approach treats them as deterministic operators. 

As a primary, real-world proof of concept, the \texttt{NDH-Surrealism} runtime translates nine distinct latent-space failure profiles directly into robotic canvas tool-paths using a physical peristaltic pump and multi-color deposition setup. However, the underlying invariance-preserving mathematics suggests application fields extending significantly beyond physical art fabrication.

% --- SECTION 2: CROSS-DOMAIN INDUSTRIAL APPLICATIONS ---
\section{Industrial Scale-Up \& Engineering Horizons}

\subsection{Fractal Architectural Topologies}
Standard architectural CAD frameworks struggle with fluidic, non-repeating structural geometries due to compounding floating-point translation errors. By mapping continuous graph Laplacians to high-dimensional spatial waves, our pipeline provides an error-clamped framework for automated 3D construction gantries depositing concrete or polymers along self-similar, structural-load vectors.

\subsection{Spectrally Optimized Solar Concentrators}
In material engineering, maximizing photon capture requires surfacing textures that mimic complex non-Euclidean geometries. By leveraging our \textit{Rotational Vortex Drive (Operator 07)} and \textit{Bilinear Mesh Leveling Matrix}, deposition robotics can coat variable surfaces with micro-structured, logarithmic spiral trajectories. This allows multi-material layers to be deposited along precise spectral angles, maximizing omnidirectional capture profiles for ultra-high-efficiency solar enclaves.

\end{document}
